%==============================================================================
%  PORTADA
%==============================================================================
\begin{titlingpage}
\begin{center}
\vspace*{2.2cm}

{\sffamily\large\color{ColorGris} TRATADO DE}

\vspace{0.4cm}
{\sffamily\HUGE\bfseries\color{ColorTinta} DERECHO CONCURSAL\\[6pt] CHILENO}

\vspace{0.9cm}
{\color{ColorBorde}\rule{0.62\textwidth}{0.8pt}}

\vspace{0.9cm}
{\large\itshape Manual sistemático de la Ley \Nro 20.720\\[4pt]
y su reforma por la Ley \Nro 21.563}

\vspace{1.6cm}
{\sffamily\normalsize\color{ColorGris}
Dogmática de la insolvencia, tramitación forense\\[3pt]
y jurisprudencia de los tribunales superiores}

\vfill

{\large\sffamily Alberto Tomás Stead Hogg}\\[4pt]
{\small\color{ColorGris} Abogado\\
Licenciado en Ciencias Jurídicas y Sociales\\
Universidad Diego Portales}

\vspace{1.4cm}
{\small\sffamily\color{ColorGris} Primera edición\\ Santiago de Chile, \the\year}

\vspace*{1.2cm}
\end{center}
\end{titlingpage}

% --- Página de créditos ---------------------------------------------------
\thispagestyle{empty}
\null\vfill
\noindent
{\footnotesize
\textbf{Tratado de Derecho Concursal Chileno.}
Manual sistemático de la Ley \Nro 20.720 y su reforma por la Ley \Nro 21.563.\par
\vspace{6pt}
\copyright\ \the\year, Alberto Tomás Stead Hogg.\par
Inscripción \Nro [XXX.XXX] en el Registro de Propiedad Intelectual.\par
ISBN [XXX-XXX-XXXX-XX-X].\par
\vspace{6pt}
Ninguna parte de esta obra puede ser reproducida sin autorización escrita del autor.\par
\vspace{10pt}
\textit{Advertencia.} Esta obra tiene finalidad doctrinaria y docente. Su contenido no
constituye asesoría jurídica para un caso concreto y no sustituye la revisión del texto
oficial de las normas ni de las resoluciones vigentes de la \Superir.\par
\vspace{6pt}
Texto cerrado al [fecha de cierre de edición]. Se recomienda contrastar el articulado
con la versión vigente en \textsc{bcn}, Ley Chile.\par
}
\vspace{1cm}
\cleardoublepage

% --- Nota sobre el autor --------------------------------------------------
\thispagestyle{empty}
\null\vspace{2cm}
\begin{center}
{\sffamily\large\bfseries\color{ColorTinta} El autor}
\end{center}
\vspace{0.6cm}
\noindent
\textbf{Alberto Tomás Stead Hogg} es abogado y Licenciado en Ciencias Jurídicas y
Sociales por la Universidad Diego Portales.\par
\vspace{8pt}
\noindent
Su formación en materia concursal proviene del ejercicio directo. Entre 2018 y 2022 se
desempeñó como procurador judicial y encargado de revisión de escritos en la liquidadora
concursal Valeria Cañas Aranda, donde intervino en la tramitación cotidiana de
procedimientos de liquidación ante los juzgados civiles de Santiago. Posteriormente fue
socio administrador de Derecho Inteligente Limitada, con especialización en derecho de
los negocios, adquisiciones concursales y asesoría a emprendimientos.\par
\vspace{8pt}
\noindent
En el ámbito académico ha estado vinculado de forma continua a la Clínica Jurídica Civil
de la Universidad Diego Portales desde 2021, primero como representante \emph{ad
honorem} y ayudante docente, luego como ayudante senior y profesor \emph{pro-tempore}
senior en el área de Innovación y Emprendimiento.\par
\vspace{8pt}
\noindent
Ha sido cofundador de plataformas de tecnología legal orientadas al seguimiento judicial
y a la automatización de flujos de trabajo, experiencia que informa el enfoque
operativo y procedimental que recorre la Parte Quinta de esta obra.\par
\vfill
\cleardoublepage
